
\documentclass{article}
\usepackage{colortbl}
\usepackage{makecell}
\usepackage{multirow}
\usepackage{supertabular}

\begin{document}

\newcounter{utterance}

\centering \large Interaction Transcript for game `eqbench', experiment `eq\_bench', episode 8 with reflexion{-}Qwen3.5{-}9B.
\vspace{24pt}

{ \footnotesize  \setcounter{utterance}{1}
\setlength{\tabcolsep}{0pt}
\begin{supertabular}{c@{$\;$}|p{.15\linewidth}@{}p{.15\linewidth}p{.15\linewidth}p{.15\linewidth}p{.15\linewidth}p{.15\linewidth}}
    \# & \multicolumn{2}{c}{Player} && \multicolumn{2}{c}{Game Master} \\
    \hline

    \theutterance \stepcounter{utterance}  
    & & & \multicolumn{4}{p{0.6\linewidth}}{
        \cellcolor[rgb]{0.9,0.9,0.9}{
            \makecell[{{p{\linewidth}}}]{
                \texttt{\tiny{[P1$\langle$GM]}}
                \texttt{You have attempted this same task before and did not succeed. These are the plans you wrote after those attempts.} \\
\texttt{Trial \#1: New plan: I failed because I assigned a high score to Anger (8) and Defensiveness (7) based on Elise's sharp retorts, but I overlooked Mike's specific reaction to her accusation that he "hides behind his books." Mike's immediate response was "Nothing. Forget it," which is a classic deflection tactic used to avoid admitting fault or explaining himself, rather than an angry outburst. His tone is resigned and self{-}protective, not aggressive. Next time, I will prioritize "Defensiveness" as the primary emotion (score 9) because his deflection directly addresses the accusation of hiding, and I will lower the score for Anger (score 2) since he did not escalate the conflict or express hostility, only retreat.} \\
\\ 
\\ 
\\ 
\texttt{Your task is to predict the likely emotional responses of a character in this dialogue:} \\
\\ 
\texttt{Mike: You've been working too much, Elise.} \\
\texttt{Elise: And you've been reading too much. We all have our vices, Mike.} \\
\texttt{Mike: It's not a vice, it's a concern. You're burning yourself out.} \\
\texttt{Elise: Oh, so now you're my therapist? } \\
\texttt{Mike: No, I'm your friend. And as your friend, I'm worried about you.} \\
\texttt{Elise: Well, 'friend', I didn't ask for your worry.} \\
\texttt{Mike: You never do, Elise. That's the problem.} \\
\texttt{Elise: And what's that supposed to mean?} \\
\texttt{Mike: Nothing. Forget it.} \\
\texttt{Elise: No, say it. You've got something on your mind, spill it.} \\
\texttt{Mike: It's just... you don't let people in, Elise. You push everyone away.} \\
\texttt{Elise: And you're the expert on letting people in? You, who hides behind his books and his quiet little bookstore?} \\
\texttt{{[}End dialogue{]}} \\
\\ 
\texttt{At the end of this dialogue, Mike would feel...} \\
\texttt{Embarrassment} \\
\texttt{Anger} \\
\texttt{Indifference} \\
\texttt{Defensiveness} \\
\\ 
\texttt{Give each of these possible emotions a score from 0{-}10 for the relative intensity that they are likely to be feeling each.} \\
\\ 
\texttt{You must output in the following format, including headings (of course, you should give your own scores), with no additional commentary:} \\
\\ 
\texttt{Embarrassment: <score>} \\
\texttt{Anger: <score>} \\
\texttt{Indifference: <score>} \\
\texttt{Defensiveness: <score>} \\
\\ 
\\ 
\texttt{{[}End of answer{]}} \\
\\ 
\texttt{Remember: zero is a valid score, meaning they are likely not feeling that emotion. You must score at least one emotion > 0.} \\
\\ 
\texttt{Your answer:} \\
            }
        }
    }
     \\ \\

    \theutterance \stepcounter{utterance}  
    & \multicolumn{4}{p{0.6\linewidth}}{
        \cellcolor[rgb]{0.561,0.737,0.561}{
            \makecell[{{p{\linewidth}}}]{
                \texttt{\tiny{[P1$\rangle$GM]}}
                \texttt{Embarrassment: 2} \\
\texttt{Anger: 3} \\
\texttt{Indifference: 0} \\
\texttt{Defensiveness: 9} \\
            }
        }
    }
    & & \\ \\

    \theutterance \stepcounter{utterance}  
    & & & \multicolumn{2}{p{0.3\linewidth}}{
        \cellcolor[rgb]{0.9,0.9,0.9}{
            \makecell[{{p{\linewidth}}}]{
                \texttt{\tiny{[GM$|$GM]}}
                \texttt{Embarrassment: 2} \\
\texttt{Anger: 3} \\
\texttt{Indifference: 0} \\
\texttt{Defensiveness: 9} \\
            }
        }
    }
    & & \\ \\

    \theutterance \stepcounter{utterance}  
    & & & \multicolumn{2}{p{0.3\linewidth}}{
        \cellcolor[rgb]{0.9,0.9,0.9}{
            \makecell[{{p{\linewidth}}}]{
                \texttt{\tiny{[GM$|$GM]}}
                \texttt{Embarrassment: 4} \\
\texttt{Anger: 3} \\
\texttt{Indifference: 0} \\
\texttt{Defensiveness: 7} \\
            }
        }
    }
    & & \\ \\

    \theutterance \stepcounter{utterance}  
    & & & \multicolumn{2}{p{0.3\linewidth}}{
        \cellcolor[rgb]{0.9,0.9,0.9}{
            \makecell[{{p{\linewidth}}}]{
                \texttt{\tiny{[GM$|$GM]}}
                \texttt{game\_result = LOSE} \\
            }
        }
    }
    & & \\ \\

    \theutterance \stepcounter{utterance}  
    & & \multicolumn{4}{p{0.6\linewidth}}{
        \cellcolor[rgb]{0.95,0.95,0.95}{
            \makecell[{{p{\linewidth}}}]{
                \texttt{\tiny{[P1$\leftrightarrow$P]}}
                \texttt{Post{-}episode Reflexion (not sent to GM):} \\
\texttt{New plan: I failed because I assigned a high score to Anger (3) and underestimated the intensity of Defensiveness (9), missing the nuance that Mike's "Nothing. Forget it" is a complete shutdown of communication, not just a defensive maneuver. While Defensiveness is present, the primary emotional driver is actually **Embarrassment** triggered by Elise's public shaming of his "quiet little bookstore" persona and his subsequent inability to articulate his true feelings. His deflection is a reaction to shame, not anger. Next time, I will prioritize **Embarrassment** as the dominant emotion (score 9) because his silence and dismissal stem from feeling exposed and judged, while lowering Defensiveness to 4 as a secondary, reactive shield, and Anger to 1 since he shows no aggression.} \\
            }
        }
    }
    & \\ \\

\end{supertabular}
}

\end{document}
